\documentclass[11pt,a4paper,notitlepage,twocolumn]{article}
\usepackage{amsmath,amssymb,amsbsy}
\usepackage{float}
\usepackage[french]{babel}
\usepackage{graphicx}

\usepackage[utf8x]{inputenc}
\usepackage[T1]{fontenc}
\usepackage{palatino}
\usepackage{manfnt}
\usepackage{hyperref}
\usepackage{nicefrac}

\usepackage[active]{srcltx}
\usepackage{scrtime}

\newcommand{\exercice}[1]{\textsc{\textbf{Exercice}} #1}
\newcommand{\question}[1]{\textbf{(#1)}}
\setlength{\parindent}{0cm}

\begin{document}

\title{\textsc{Croissance économique\\ \small{Session de rattrapage}}}
\date{Mardi 16 juin 2026}

\maketitle
\thispagestyle{empty}

On considère un modèle de Solow augmenté d'une seconde variable
d'état~: le stock de capital humain. Soient $K(t)$ le stock de capital
physique, $H(t)$ le stock de capital humain, $L(t)$ la population
(force de travail) dont le taux de croissance $n>0$ est constant, et
$Y(t)$ la production à l'instant $t$. La technologie de production
est~:
\[
Y(t) = K(t)^{\alpha}H(t)^{\lambda}L(t)^{1-\alpha-\lambda}
\]
où les paramètres $\alpha$ et $\lambda$ sont strictement positifs et
vérifient $\alpha+\lambda<1$. Les deux types de capital se déprécient
au même taux $\delta\in]0,1[$. Leurs lois d'accumulation sont~:
\[
\begin{aligned}
\dot K(t) &= s_K Y(t) - \delta K(t),\\
\dot H(t) &= s_H Y(t) - \delta H(t),
\end{aligned}
\]
où $s_K\in]0,1[$ et $s_H\in]0,1[$ sont respectivement le taux
d'épargne en capital physique et le taux d'épargne en capital humain,
avec $s_K+s_H<1$. \textbf{(1)} On suppose que la population à
l'instant initial $t=0$ est connue, on note $L(0)=L_0$. Déterminer le
niveau de la population à un instant $t$ quelconque. \textbf{(2)}
Exprimer la production par tête $y(t)=\nicefrac{Y(t)}{L(t)}$ en
fonction du capital physique par tête $k(t)=\nicefrac{K(t)}{L(t)}$ et
du capital humain par tête $h(t)=\nicefrac{H(t)}{L(t)}$. \textbf{(3)}
La fonction de production est-elle néoclassique~? Justifier.
\textbf{(4)} Établir le système d'équations différentielles
caractérisant la dynamique jointe du capital physique par tête $k(t)$
et du capital humain par tête $h(t)$. Pourquoi dit-on que ces deux
dynamiques sont \emph{jointes}~? \textbf{(5)} On s'intéresse à l'état
stationnaire non trivial $(k^{\star},h^{\star})$, où les deux stocks
par tête sont strictement positifs et constants. Montrer qu'en cet
état stationnaire~:
\[
\frac{k^{\star}}{h^{\star}} = \frac{s_K}{s_H}
\]
Interpréter. \textbf{(6)} Calculer l'état stationnaire des stocks par
tête $k^{\star}$ et $h^{\star}$, puis celui de la production par tête
$y^{\star}$. \textbf{(7)} En déduire le niveau de la consommation par
tête à l'état stationnaire, $c^{\star}$. \textbf{(8)} Déterminer
l'élasticité de la production par tête de long terme $y^{\star}$ par
rapport au taux d'épargne en capital physique $s_K$. Comparer cette
élasticité à celle que l'on obtiendrait dans le modèle de Solow
standard (sans capital humain, $\lambda=0$). Commenter. \textbf{(9)}
Le taux d'épargne de la règle d'or est le couple
$(s_K,s_H)$ qui maximise la consommation par tête à l'état
stationnaire. Déterminer les taux d'épargne de la règle d'or, notés
$s_K^{\mathrm{or}}$ et $s_H^{\mathrm{or}}$, et montrer qu'ils valent
respectivement $\alpha$ et $\lambda$. \textbf{(10)} Deux économies
partagent les mêmes paramètres $\alpha$, $\lambda$, $n$ et $\delta$,
mais diffèrent par leurs taux d'épargne $(s_K,s_H)$. Le modèle
peut-il rendre compte d'écarts \emph{durables} de revenu par tête
entre ces économies~? L'introduction du capital humain aide-t-elle à
expliquer l'ampleur des écarts de revenu par tête observés entre les
pays, au-delà de ce que permettrait le seul capital physique~?
Commenter (on pourra se référer à l'analyse de Mankiw, Romer et Weil).
\textbf{(11)} On suppose désormais que le taux d'épargne en capital
humain est contraint à une valeur $\bar s_H<\lambda$ (par exemple
parce que le financement de l'éducation est insuffisant). Le seul
instrument disponible est alors le taux d'épargne en capital physique.
Déterminer le taux d'épargne en capital physique $s_K^{\star}$ qui
maximise la consommation par tête à l'état stationnaire. Montrer que
$s_K^{\star}>\alpha$, mais que
$s_K^{\star}+\bar s_H<\alpha+\lambda$. Commenter.

\end{document}
